\section{Institutional background}
\label{sec:institutions}

The statutory minimum wage was created by Decreto-Lei n.º 217/74 of 27 May 1974
and has been set by decree ever since, normally with effect from 1 January.
Three features of the institution matter for measurement and are routinely lost
in secondary compilations.

\paragraph{The wage was not single until 2005.} The law distinguished separate
regimes, and the number fell over time. Until 1991 three coexisted: domestic
service, agriculture, and a general regime covering the remaining activities.
From 1992 agriculture was folded into the general regime, leaving two. Only from
2005 is there a single national figure. The general regime is the series
conventionally quoted as \emph{the} Portuguese minimum wage and is the one used
throughout; the special regimes were always set at or below it.

\paragraph{Payment is fourteen times a year.} The statutory figure is monthly,
and the law provides for holiday and Christmas subsidies, so annual statutory
pay is fourteen times the monthly value rather than twelve. International
compilations often express the wage on a twelve-month basis, which produces a
figure about seventeen per cent higher than the one in Portuguese law and is a
common source of confusion.

\paragraph{The autonomous regions may set their own wage, by different
mechanisms.} This distinction is central to
Section~\ref{sec:robustness}. The Azores legislate a permanent proportional
supplement: article~3 of Decreto Legislativo Regional n.º 8/2002/A adds five per
cent to whatever the national wage is. Madeira instead legislates a value each
year, and its premium over the mainland is not a rule: it has widened from
about three per cent in 2023 to over six per cent in 2026.

The distinction is not administrative. A proportional supplement is a fixed
transformation of the national wage, so in logs the Azorean statutory change
equals the mainland's in every month after the supplement was introduced.
Madeira's does not. Whatever regional identification is available in Portugal
comes from Madeira and from the single month in which the Azorean supplement
took effect.

\paragraph{Escudo and euro.} Values before 2002 are in escudos and convert at
200.482 to the euro, the irrevocable rate fixed for the changeover. The act
effective in January 2002 states both currencies, and its published euro figure
is used in preference to the converted one.
